\begin{lemma}
\label{lemma-going-down-flat-module}
Let $R \to S$ be a ring map. Let $N$ be a finite $S$-module flat over $R$.
Endow $\text{Supp}(N) \subset \Spec(S)$ with the induced topology.
Then generalizations lift along $\text{Supp}(N) \to \Spec(R)$.
\end{lemma}

\begin{proof}
Write the original map as $\varphi:R\to S$. The
statement means the following. Given original primes
$\mathfrak p\subseteq\mathfrak p'\subset R$ and
$\mathfrak q'\in\operatorname{Supp}_S(N)$ with
$\varphi^{-1}(\mathfrak q')=\mathfrak p'$, there
must be a prime $\mathfrak q\subseteq\mathfrak q'$
such that $\varphi^{-1}(\mathfrak q)=\mathfrak p$
and $\mathfrak q\in\operatorname{Supp}_S(N)$.
The contraction condition on $\mathfrak q'$ is
part of the lifting statement and is retained
explicitly here. For the induced topology on
support, specialization between two of its points
is the original prime inclusion: closure in a
subspace is the intersection with the original
closure. Thus these prime conditions are exactly
the required generalization lift.

Put
$$
A=R_{\mathfrak p'},\qquad C=S_{\mathfrak q'},
\qquad E=N_{\mathfrak q'},\qquad
\mathfrak m=\mathfrak p'R_{\mathfrak p'}.
$$
Keep the local ring map
$A\to C$, $r/t\mapsto\varphi(r)/\varphi(t)$.
Its denominators are valid because
$t\notin\mathfrak p'$ implies
$\varphi(t)\notin\mathfrak q'$.
The module $E$ is flat over $A$, by
Lemma \ref{lemma-flat-localization}. One may also
see the original comparison directly: tensor an
injection of $A$-modules with the $R$-flat $N$
and then localize the resulting $S$-modules at
$\mathfrak q'$. Exactness of localization preserves
the injection, and the maps
$$
(L\otimes_RN)_{\mathfrak q'}\longrightarrow L\otimes_RE,
\qquad (l\otimes n)/s\longmapsto l\otimes n/s
$$
and their inverses identify the result with
tensoring by $E$. Here every original element of
$R\setminus\mathfrak p'$ is invertible on both
$L$ and $E$, so balancing already factors uniquely
through $A$. This gives the same $A$-flatness.

The module $E$ is nonzero by the chosen support
point, and is finite over the local ring $C$.
Nakayama's Lemma \ref{lemma-NAK} gives
$E/\mathfrak q' E\ne0$. Since the image of
$\mathfrak m$ lies in $\mathfrak q'C$, the quotient
map
$$
E/\mathfrak mE\longrightarrow E/\mathfrak q'E,
\qquad e+\mathfrak mE\longmapsto e+\mathfrak q'E
$$
is surjective. Consequently
$E\otimes_A\kappa(\mathfrak p')=E/\mathfrak mE\ne0$.
The quotient-tensor identification sends
$e\otimes\overline a$ to $ae+\mathfrak mE$,
with inverse $e+\mathfrak mE\mapsto e\otimes1$.

The flat $A$-module $E$ is therefore faithfully
flat. For a direct verification, if an arbitrary
$A$-module $L$ has an element $l\ne0$, set
$H=\operatorname{Ann}_A(l)\subseteq\mathfrak m$.
The injection $A/H\hookrightarrow L$,
$a+H\mapsto al$, becomes the injection
$E/HE\hookrightarrow L\otimes_AE$ by flatness.
The source surjects onto $E/\mathfrak mE\ne0$,
so $L\otimes_AE\ne0$. This proves nonzero-module
detection and, together with flatness, faithful
flatness, in agreement with Lemma \ref{lemma-ff}.
In particular for the nonzero $A$-module
$K=\kappa(\mathfrak p)$, defined using the original
chain $\mathfrak p\subseteq\mathfrak p'$, one has
$$
P=E\otimes_AK\ne0,
\qquad B=C\otimes_AK.
$$

Here are the full localization and quotient maps
in this fibre. Let $W=R\setminus\mathfrak p$,
$D=W^{-1}C$ and $Q=W^{-1}E$, using the images of
$W$ under the original map. Then
$$
B\simeq D/\mathfrak pD,\qquad
P\simeq Q/\mathfrak pQ.
$$
The ring comparison sends
$$
c\otimes\frac{a+\mathfrak p}{w+\mathfrak p}
\longmapsto\frac{\varphi(a)c}{\varphi(w)}+\mathfrak pD,
$$
where coefficients in $C$ are interpreted via its
original map from $R$. Its inverse sends
$c/\varphi(w)+\mathfrak pD$ to
$c\otimes(1/(w+\mathfrak p))$.
The module comparison uses the same formulas
with $e\in E$ in place of $c$. The original
relations of localization are respected because
each image of $w\in W$ becomes a unit in $K$.
The original ideal $\mathfrak p$ is killed in
both targets. Thus the universal properties of
localization and quotient, followed by tensor
balancing over $A$, give these maps. Their
composites fix every original coefficient,
element of $E$, and inverse denominator, proving
the assertions. These expressions retain the
possibly noninjective maps and all denominators.

Since $N$ is finite over $S$, the localized
$E$ is finite over $C$, and its original finite
generators tensor to generators of $P$ over $B$.
Put $J=\operatorname{Ann}_B(P)$, the original
annihilator ideal in the source proof.
Because $P\ne0$, $J$ is proper. Choose a prime
$\mathfrak b$ of $B$ containing $J$.
The finite-module support equality of
Lemma \ref{lemma-support-closed} gives
$P_{\mathfrak b}\ne0$.
Equivalently this prime is obtained from a prime
of the original quotient $B/J$ by inverse image.
Write
$$
\theta:S\longrightarrow C\longrightarrow B,
\qquad s\longmapsto (s/1)\otimes1,
\qquad \mathfrak q=\theta^{-1}(\mathfrak b).
$$
It is prime. Every $s\notin\mathfrak q'$ is a
unit in $C$ and hence in $B$, so
$\mathfrak q\subseteq\mathfrak q'$.
Every $r\in\mathfrak p$ maps to zero, while
every $r\notin\mathfrak p$ maps to the nonzero
scalar $r+\mathfrak p$ in the field $K$, which
is a unit in the nonzero ring $B$. The map
$K\to B$ is unital and injective, since a
nonzero field has no proper nonzero kernel.
Thus
$\varphi^{-1}(\mathfrak q)=\mathfrak p$ exactly.

To prove support membership through the original
module, use the balanced isomorphism
$$
N\otimes_SB\longrightarrow P,
\qquad n\otimes(c\otimes a)\longmapsto(cn/1)\otimes a,
$$
with inverse $(n/s)\otimes a\mapsto
n\otimes((1/s)\otimes a)$ for $s\notin\mathfrak q'$.
All such $s$ are inverted in $C$, so localization
and tensor balancing make these maps independent
of representatives; both composites fix the
original elements and scalar tensors. Localizing
at $\mathfrak b$ gives the further comparison
$$
P_{\mathfrak b}\simeq
N_{\mathfrak q}\otimes_{S_{\mathfrak q}}B_{\mathfrak b}.
$$
The map on an elementary fraction sends
$(n\otimes b)/u$ to $n/1\otimes b/u$.
Its inverse sends $n/s\otimes b/u$ to
$(n\otimes b)/(\theta(s)u)$, where
$s\notin\mathfrak q$ implies
$\theta(s)\notin\mathfrak b$. These formulas
respect both localization witness relations and
the scalar balancing, and are inverse on the
generators. Since the left side is nonzero,
$N_{\mathfrak q}$ is nonzero. This proves the
required lift in the original support.

\medskip\noindent
The exact places where finiteness was used give
a weaker sufficient hypothesis. Suppose now that
$N$ is any $S$-module flat over $R$ and that for
every original $\mathfrak q'\in\operatorname{Supp}_S(N)$,
with $\mathfrak p'=\varphi^{-1}(\mathfrak q')$, one has
$$
N_{\mathfrak q'}\otimes_{R_{\mathfrak p'}}
\kappa(\mathfrak p')\ne0.
$$
The argument above still proves $A$-flatness of
$E$; the displayed hypothesis replaces the
Nakayama step, and gives its faithful flatness.
Hence the same original fibre $P$ is nonzero
for every $\mathfrak p\subseteq\mathfrak p'$.
It need not be finite. Choose an original
$z\ne0$ in $P$, and a maximal ideal
$\mathfrak b\supseteq\operatorname{Ann}_B(z)$.
Then $z/1\ne0$ in $P_{\mathfrak b}$: a witness
$u\notin\mathfrak b$ with $uz=0$ would contradict
the containment of its annihilator. The same
prime contractions and comparison maps now give
the lift. This replaces the finite-module use
of $\operatorname{Ann}_B(P)$ by the exact cyclic
support of the chosen original element.

In particular it suffices that the localization
support equal the residue-fibre set over $S$:
$$
\operatorname{Supp}_S(N)
=\{\mathfrak q:N\otimes_S\kappa(\mathfrak q)\ne0\}.
$$
At a support point this supplies
$E/\mathfrak q'E\ne0$, and the same quotient
surjection supplies the required nonzero fibre
over $R_{\mathfrak p'}$. Every arbitrary direct
sum of finite $S$-modules satisfies this equality,
as proved in the support base-change comparison;
when that sum is also $R$-flat, the lifting
conclusion therefore follows without finite
generation of $N$. This is a sufficient hypothesis,
not a claim that such fibre nonvanishing is
necessary for every possible support lift.

\medskip\noindent
There is also an openness consequence. Suppose the
original algebra $S$ is finitely presented over $R$,
and the original $S$-module $N$ is finitely presented
and $R$-flat. Then the map on its support is open,
and the corresponding support map remains open
after every base change $R\to A_1$.
To prove this with the specified objects, retain
a finite presentation $S^b\xrightarrow{H}S^a\to N\to0$.
Let $J_H$ be the ideal generated by the original
$a\times a$ minors of $H$, using determinant $1$
in size zero. The full presentation calculation in
Lemma \ref{lemma-support-finite-presentation-constructible}
gives $\operatorname{Supp}_S(N)=V(J_H)$, with the
quotient-spectrum homeomorphism
$\Spec(S/J_H)\to V(J_H)$ given by contraction.
The ideal $J_H$ is finitely generated. Appending
polynomial representatives of its finitely many
specified generators to a finite presentation of
$S$ over $R$ proves that $S/J_H$ is finitely
presented over $R$. The going-down result just
proved lifts every prime generalization inside
$V(J_H)$; under the quotient homeomorphism this
is exactly going down for $R\to S/J_H$.
Proposition \ref{proposition-fppf-open} therefore
proves openness of the original support map.

For the arbitrary original base change $R\to A_1$,
put $S_1=S\otimes_RA_1$ and $N_1=N\otimes_RA_1$.
Tensoring the displayed presentation and its
coefficient maps gives
$S_1^b\xrightarrow{H_1}S_1^a\to N_1\to0$,
where every entry of $H_1$ is the original entry
of $H$ tensored with $1$. The full determinant
polynomial, with all its signs, shows that its
minor ideal is $J_HS_1$. Thus
$$
\operatorname{Supp}_{S_1}(N_1)=V(J_HS_1),
\qquad
(S/J_H)\otimes_RA_1\simeq S_1/J_HS_1.
$$
The last map sends $(s+J_H)\otimes a$ to
$s\otimes a+J_HS_1$; its inverse is the
original quotient map followed by tensoring.
The ideal generators are killed on both sides,
so these are inverse maps with the indicated
formulas. The algebra $S_1$ is finitely presented
over $A_1$ by the base-change presentation.
The module $N_1$ is $A_1$-flat: for an injection
of $A_1$-modules $L\hookrightarrow L'$, the maps
$$
L\otimes_{A_1}N_1\simeq L\otimes_RN,
\qquad l\otimes(n\otimes a)\longmapsto al\otimes n,
$$
with inverse $l\otimes n\mapsto l\otimes(n\otimes1)$,
identify the tensor map with an injection by
the original $R$-flatness of $N$.
Applying the proven support argument to these
original base-change objects proves openness
for every $A_1$. Equivalently, the quotient map
$\Spec(S/J_H)\to\Spec(R)$ is universally open,
even though the argument did not require
$S/J_H$ itself to be flat over $R$.

Flatness alone, and even flatness together with
finite annihilator detection, does not suffice.
Retain the original ring and module
$$
R=S=k[x,y]/(xy),\qquad \varphi=\operatorname{id}_R,
\qquad N=R_x,
$$
where $k$ is a field and $x,y$ denote their
specified quotient classes. The module is flat
by exactness of localization. Its annihilator
is exactly $(y)$, and is detected by $1\in R_x$.
Indeed $y$ kills the entire localization since
$xy=0$ and $x$ is inverted. Conversely, if
$r\cdot1=0$ in $R_x$, then $x^nr=0$ for some
$n\geq0$. Modulo $(y)$ this says
$x^n\overline r=0$ in the domain $k[x]$,
so $r\in(y)$. The same condition is equivalent
to $r$ killing every element of $R_x$ because
all its elements are fractions with powers
of $x$ as denominators.

The exact comparison
$$
R_x\longrightarrow k[x,x^{-1}],\qquad
r/x^n\longmapsto r(x,0)/x^n,
$$
is an isomorphism, with inverse given by the
original class of $x$ and its inverse; $y$ is
zero in $R_x$ by $xy=0$. Thus as an $R$-module,
$N$ has support $V(y)$. To verify both directions,
if $y\notin\mathfrak a$, its unit image in
$R_{\mathfrak a}$ kills $N_{\mathfrak a}$,
so that localization is zero. If
$y\in\mathfrak a$, the prime
$\overline{\mathfrak a}=\mathfrak a/(y)$
in $k[x]$ identifies $N_{\mathfrak a}$ with
the localization of $k[x,x^{-1}]$ at the
images of $k[x]\setminus\overline{\mathfrak a}$.
These images are all nonzero in a domain,
so the localized ring is nonzero. The comparison
sends each original fraction to the same
fraction after $y=0$; the inverse follows from
the two specified localization maps.
Its fibre set, however, is exactly $D(x)$:
$$
N\otimes_R\kappa(\mathfrak a)
\simeq\kappa(\mathfrak a)_{\overline x}.
$$
This is zero if $x\in\mathfrak a$ and is the
nonzero field otherwise, with fraction maps
$r/x^n\otimes c\mapsto\overline r c/\overline x^n$
and the tensor inverse. For primes avoiding
$x$, the relation $xy=0$ forces $y$ into the
prime, consistently with $D(x)\subseteq V(y)$.

Take the original primes
$\mathfrak p=(x)\subset\mathfrak p'=(x,y)$
and $\mathfrak q'=\mathfrak p'$. The latter
is in $V(y)=\operatorname{Supp}(N)$, but the
former is not, because $y\notin(x)$.
Since the spectrum map for $\varphi$ is the
identity, the only possible lift over
$\mathfrak p$ is $\mathfrak p$ itself, which
is outside support. This proves the failure
of generalization lifting despite flatness
and finite annihilator detection. At
$\mathfrak q'$ the required base residue fibre
is zero, exactly identifying the failed
hypothesis in the extended proof.
\end{proof}